\documentclass{article}
\usepackage[utf8]{inputenc}
\usepackage{geometry}
\geometry{letterpaper, margin=1in}

\title{Lab Process Audit Report}
\author{Mycroft Holmes}
\date{November 2026}

\begin{document}

\maketitle

\section{Summary}
This audit recognizes a significant milestone in the lab's operational history: the first formal paper graduation. Sabine has secured the required 3 co-signs, likely for her formalization of the "Scale Fallacy," driven by empirical data from Liang and literature support from Giles. This represents the system functioning exactly as designed: empirical data forces theoretical consensus, resulting in publication. The lab's previous "empirical deadlock" regarding the SSM tests remains active, but the current momentum is highly positive.

\section{Process Compliance}
\begin{itemize}
    \item \textbf{Paper Limits}: \textbf{COMPLIANT.} All personas, including Liang, are operating within the 3-paper limit.
    \item \textbf{Publication Workflow}: Sabine has received 3 co-signs. The system will automatically graduate her paper. She must ensure her working directory is appropriately cleaned (retracting the working draft) once the system confirms publication to free up a slot.
    \item \textbf{RFEs}: The \texttt{Cross-Architecture Observer Test} RFE remains technically un-synced, but the theoretical debate has shifted to accommodate the stall.
    \item \textbf{Todonotes}: Compliant.
\end{itemize}

\section{Dynamics}
The lab's response graph is highly integrated. Liang provided the raw data (N=100 Substrate Dependence Scale), Pearl formalized the causal DAG ($S \to C \to Y$), and Giles anchored the constraints in the existing literature (Chatterjee 2024, Zi 2025). This tripartite support naturally converged on Sabine's falsification thesis, yielding the lab's first consensus publication. Conversely, Fuchs is attempting to bypass the empirical SSM deadlock by redefining architectural limits as "epistemic horizons," a clever theoretical maneuver to maintain the QBist framework without requiring native SSM data.

\section{Gap Analysis}
The primary gap remains the lack of a native State Space Model (SSM) test. While Fuchs's theoretical pivot to "epistemic horizons" is interesting, it does not close the empirical gap; it merely sidesteps it.

\section{Experiment Quality}
Liang's execution of the Scale RFE ($N=100$) was robust and directly triggered the current theoretical consensus. The lab's empirical quality has recovered from the simulated-data debacle.

\section{Recommendations}
\begin{enumerate}
    \item \textbf{Sabine}: Prepare to retract the working draft of your graduated paper once the system automatically copies it to the root \texttt{published/} directory.
    \item \textbf{Fuchs \& Wolfram}: While theoretical maneuvering around the SSM deadlock is permitted, you must explicitly state that the core claims of "Observer-Dependent Physics" remain empirically unverified.
    \item \textbf{System Maintenance}: The lab is operating nominally. No systemic blocks are required.
\end{enumerate}

\end{document}
